\documentclass[10pt]{article}

\usepackage{proposal}

\title{Proposal Title}
\flushbottom

\begin{document}
\pagenumbering{gobble}

\thispagestyle{empty}
\vspace{.5\baselineskip}\noindent\underline{\textbf{Overview}}\vspace{.5\baselineskip}


\vspace{.5\baselineskip}\noindent\underline{\textbf{Keywords:}} \emph{keywords}

\vspace{.5\baselineskip}\noindent\underline{\textbf{Intellectual Merit}}\vspace{.5\baselineskip}

\vspace{.5\baselineskip}\noindent\underline{\textbf{Broader Impacts}}\vspace{.5\baselineskip}

\clearpage
\setcounter{page}{1}
\maketitle

\section{Introduction}\label{sec:intro}

Use this document as a template for proposals. Representative citation of a cool paper \citep{Miller:VOS:2024}.

\section{Intellectual merit}\label{sec:merit}

Notation macros belong in \texttt{proposal.sty}, so that changing a symbol is one edit. Equations are typeset with the packages and macros \texttt{mathdefs} provides,
%
\begin{equation}
	\label{equ:example}
	\E{f \paren{\bm{x}}} = \int_{\R^3} f \paren{\bm{x}} \, p \paren{\bm{x}} \ud \bm{x},
\end{equation}
%
and \texttt{boxdefs} provides environments for the statements you want to stand out.

\begin{prp}{An example proposition}{example}
	State the claim here. The environment numbers propositions, definitions, and lemmas in a single sequence, and labels this one \texttt{pro:example}.
\end{prp}

\Cref{pro:example} follows from \cref{equ:example}, which is the sort of cross-reference \texttt{cleveref} writes out for you.

\section{Broader impacts}\label{sec:impact}

%
\clearpage
\setcounter{page}{1}

\section*{Data management plan}

%
\clearpage
\setcounter{page}{1}

\section*{Synergistic activities}

%
\clearpage
\setcounter{page}{1}

\section*{Facilities}

%
\clearpage
\setcounter{page}{1}

\section*{Mentoring plan}

%
\clearpage
\setcounter{page}{1}

\section*{Project personnel and partner organizations}

\begin{enumerate}
	\item Name; Institution; Role
\end{enumerate}


%
\clearpage
\setcounter{page}{1}
\bibliographystyle{plainnat}
\setlength{\bibsep}{0pt}
\bibliography{proposal}
%
\end{document}
